\begin{abstract}
A phishing URL detector is trained on one collection but deployed on traffic it has never seen. We ask whether, when \emph{no} target data are available for fitting, feature selection or threshold calibration, choosing lexical features by the \emph{stability} of their source-domain importance transfers better than choosing them by mean importance at the same feature budget, and what such compression costs. EBBU-2017 and HISPAR--Phishstats (HP) form a curated, domain-disjoint benchmark evaluated in both directions with random forest and LightGBM, nested source-only budget selection, source-locked thresholds and a frozen protocol. Transfer is asymmetric: EBBU-trained LightGBM reaches F1 \res{F1ABLGBMALL} on HP, whereas HP-trained models reach \res{F1BARFALL}--\res{F1BALGBMALL} on EBBU. At matched budgets, stability ranking improves HP$\rightarrow$EBBU F1 by \res{DeltaBARF} (RF) and \res{DeltaBALGBM} (LightGBM) and is neutral in the other direction, yet no reduced subset matches all 37 features. Exploratory diagnostics on validation data explain why: a logistic-regression baseline keeps ranking ability but collapses at the source threshold, a single ``not HTTPS'' rule already reaches F1 \res{xRuleTgtA} on HP, and \res{xNFlip} of \res{xNEda} inspected features reverse their class association between datasets. Stability measured inside one source cannot see such reversals. The contribution is an auditable empirical comparison with a bounded insight, not a new learning algorithm or an improvement over domain adaptation.
\end{abstract}
